% ===== PRÉAMBULE CANONIQUE — à coller VERBATIM en tête de CHAQUE .tex =====
\documentclass[11pt,a4paper]{article}
\usepackage[utf8]{inputenc}\usepackage[T1]{fontenc}\usepackage{lmodern}
\usepackage[french]{babel}
\usepackage{amsmath,amssymb,mathtools}\usepackage{icomma}
\usepackage[version=4]{mhchem}          % \ce{...} : équations & formules
\usepackage{siunitx}\sisetup{locale=FR} % \qty{}{}, \si{}, \ang{}
\usepackage{chemfig}                    % \chemfig{...} : structures développées
\usepackage{xcolor}\usepackage{colortbl}
\usepackage{tikz}\usepackage{pgfplots}\pgfplotsset{compat=1.18}
\usepackage[most]{tcolorbox}\usepackage{enumitem}\usepackage{array,booktabs}
\usepackage[a4paper,margin=2.2cm]{geometry}
\usetikzlibrary{arrows.meta,calc,angles,quotes,positioning,patterns,backgrounds,shapes.geometric}
\definecolor{primary}{RGB}{222,49,99}\definecolor{primarydark}{RGB}{150,22,55}
\definecolor{defcol}{RGB}{0,114,178}\definecolor{methcol}{RGB}{52,92,148}
\definecolor{excol}{RGB}{0,135,90}\definecolor{attncol}{RGB}{200,40,55}
\tcbset{boxbase/.style={enhanced,breakable,boxrule=0.9pt,arc=2.5pt,
  left=3mm,right=3mm,top=2.3mm,bottom=2.3mm,fonttitle=\bfseries,coltitle=white}}
% --- boîtes TUTORIEL (noms EXACTS, ne pas renommer) ---
\newtcolorbox{cle}[1][]{boxbase,colback=defcol!6,colframe=defcol,
  title={\textsc{À retenir}\ifx\relax#1\relax\relax\else\textmd{ : \itshape #1}\fi}}
\newtcolorbox{methode}[1][]{boxbase,colback=methcol!7,colframe=methcol,sharp corners,
  title={\textsc{Méthode}\ifx\relax#1\relax\relax\else\textmd{ --- \itshape #1}\fi}}
\newtcolorbox{exemplebox}[1][]{boxbase,colback=excol!5,colframe=excol,
  title={\textsc{Exemple}\ifx\relax#1\relax\relax\else\textmd{ : \itshape #1}\fi}}
\newtcolorbox{piege}[1][]{boxbase,colback=attncol!6,colframe=attncol,
  title={\textsc{Piège}\ifx\relax#1\relax\relax\else\textmd{ : \itshape #1}\fi}}
\newtcolorbox{remarque}[1][]{boxbase,colback=black!4,colframe=black!45,
  title={\textsc{Remarque}\ifx\relax#1\relax\relax\else\textmd{ : \itshape #1}\fi}}
% --- boîtes EXERCICE / CORRIGÉ (noms EXACTS, auto-comptées) ---
\newtcolorbox[auto counter]{exo}[1][]{boxbase,colback=primary!4,colframe=primary,
  title={\textsc{Exercice}~\thetcbcounter\ifx\relax#1\relax\relax\else\textmd{ --- \itshape #1}\fi}}
\newtcolorbox[auto counter]{corrige}[1][]{boxbase,colback=excol!4,colframe=excol,
  title={\textsc{Corrigé}~\thetcbcounter\ifx\relax#1\relax\relax\else\textmd{ --- \itshape #1}\fi}}
\newcommand{\R}{\mathbb{R}}\newcommand{\N}{\mathbb{N}}\newcommand{\Z}{\mathbb{Z}}\newcommand{\ds}{\displaystyle}
% (un \usepackage{fancyhdr}+en-tête est possible ; il est ignoré côté web)
% =========================================================================
\begin{document}
\begin{exo}[l'ion formé et son gaz noble]
Pour chaque atome, donne l'ion formé et le gaz noble dont il adopte la configuration.
\textit{Données :} $Z(\ce{O})=8$, $Z(\ce{Na})=11$, $Z(\ce{Mg})=12$, $Z(\ce{K})=19$ ; gaz nobles
$\ce{Ne}$ ($10$), $\ce{Ar}$ ($18$).
\begin{enumerate}[label=\alph*)]
  \item $\ce{O}$
  \item $\ce{Na}$
  \item $\ce{Mg}$
  \item $\ce{K}$
\end{enumerate}
\end{exo}

\begin{exo}[quels ions ont la configuration d'un gaz noble ?]
Pour chaque ion, donne son nombre d'électrons et dis s'il a la configuration d'un gaz noble (lequel).
\textit{Données :} $Z(\ce{H})=1$, $Z(\ce{F})=9$, $Z(\ce{Na})=11$, $Z(\ce{Ca})=20$ ; gaz nobles
$\ce{He}$ ($2$), $\ce{Ne}$ ($10$), $\ce{Ar}$ ($18$).
\begin{enumerate}[label=\alph*)]
  \item $\ce{F^{-}}$
  \item $\ce{Na^{+}}$
  \item $\ce{Ca^{2+}}$
  \item $\ce{H^{+}}$
\end{enumerate}
\end{exo}

\begin{exo}[famille à partir du numéro atomique]
Établis la configuration, compte les électrons de valence, puis donne la famille de chaque élément.
\begin{enumerate}[label=\alph*)]
  \item $Z = 11$
  \item $Z = 17$
  \item $Z = 18$
  \item $Z = 20$
\end{enumerate}
\end{exo}

\begin{exo}[classer par électronégativité]
Classe les éléments par électronégativité \textbf{croissante} (elle augmente vers la droite d'une période).
\begin{enumerate}[label=\alph*)]
  \item période 3 : $\ce{Na}$, $\ce{Al}$, $\ce{P}$, $\ce{Cl}$ ($Z = 11, 13, 15, 17$) ;
  \item période 2 : $\ce{Li}$, $\ce{N}$, $\ce{F}$ ($Z = 3, 7, 9$).
\end{enumerate}
\end{exo}

\begin{exo}[classer par rayon atomique]
Classe les éléments par rayon atomique \textbf{croissant}.
\begin{enumerate}[label=\alph*)]
  \item même colonne (alcalins) : $\ce{Li}$, $\ce{Na}$, $\ce{K}$ ;
  \item même période (période 3) : $\ce{Na}$, $\ce{Mg}$, $\ce{Cl}$.
\end{enumerate}
\end{exo}

\end{document}
